\section{Discussion and conclusion}
\label{sec:discussion}

The central inference problem is asymmetric. A successful decoder demonstrates surviving information within its evaluation regime; a failed decoder cannot exclude every alternative code. A counterexample disproves a universal implication from contraction or low dimension to forgetting; it does not establish how often a trained model uses that counterexample. Our analytic controls exploit this asymmetry to delimit permissible explanations.

The affine null is especially useful: equal-horizon cue propagation is unchanged by body order, although offsets can alter native error. In nonlinear networks, input-dependent gating, changing operating points, and basin transitions remove this restriction, but observing such ingredients does not identify which causes a failure. Finite cue separation, task-sensitive propagation, and controlled interventions are needed together. Likewise, lower dimension with preserved performance is compatible with sufficient compression; contracting local spectra with separated cue basins are compatible with categorical memory.

The pilots in \cref{fig:pilots} expose an earlier failure of inference. Failed native competence makes later poor performance an ambiguous retention measurement; it does not localize failure to acquisition rather than query/readout. The resulting contribution is a falsification framework supported by analytic constructions and an empirical audit. Whether matched repetition erases acquired information in competent trained models remains an open empirical question.
